
Models achieving high accuracy on standard benchmarks often exhibit substantial performance degradation on held-out test sets. ImageNet classifiers drop 11-15\% on ImageNet-V2, and BERT's 84\% MNLI accuracy collapses on HANS. However, no predictive metric exists to identify saturated benchmarks before investing in leaderboard optimization. This paper proposes the Difficulty-Normalized Saturation Index (DNSI), which measures benchmark saturation via the entropy of improvement patterns in state-of-the-art histories, normalized by task difficulty. The underlying observation is that saturated benchmarks exhibit compressed improvement entropy: remaining gains are narrow, incremental optimizations rather than diverse innovations. In a pilot study across four benchmarks with published generalization gaps, DNSI correlates with gap magnitude (Pearson R = -0.950, p = 0.050, n = 4). Pre-2019 DNSI values show correlation with post-2019 generalization gaps ($R^2$ = 0.349), and the correlation direction is consistent across vision (R = -0.972) and NLP (R = -0.684) domains. These results are preliminary given the small sample size; bootstrap confidence intervals span the full range due to n = 4. DNSI can be computed from historical SOTA records without constructing held-out test sets.
